\chapter{The exact local signed EBU quote}\label{ch:quote}

\section{The action finally has a number}
We have taken care to identify a physical state, a reference, a potential and
an action that can actually occur. Only now is it safe to attach a number to
the action. The number is not the quantity delivered, the price paid or the
initial slope of the potential. It is the complete change in the declared
potential between the accepted predecessor and successor.

Let $x^{\mathrm{b}}$ be the state immediately before an action, let $q$ be its finite
source-side quantity, and let $S_e$ be the declared increment per unit of
that quantity. Under one fixed field $V$,
\begin{equation}\label{eq:exact-local-ebu}
 E_e(q)=V(x^{\mathrm{b}})-V(x^{\mathrm{b}}+S_e q).
\end{equation}
Positive $E_e$ means the represented potential fell. Negative $E_e$ means it
rose. Zero means the two endpoints have equal potential; it does not prove
that the action was empty. The complete increment and the version of $V$
belong beside the number on any trustworthy quote.

\begin{lawbox}{The finite definition}
For a fixed, declared potential, exact finite EBU is before-potential minus
after-potential for the complete accepted physical change. It is not defined
by the local gradient approximation, a monetary price or an extra charge
invented after the event.
\end{lawbox}

\section{Why a local computation can give the whole answer}
Suppose the potential is a sum of factors, $V=\sum_\alpha\phi_\alpha$, each
with a declared set of input coordinates. An action changes only some of
those coordinates. In $V(x^{\mathrm{b}})-V(x^{\mathrm{b}}+S_eq)$ every factor whose inputs are unchanged
cancels exactly. Thus a local evaluator needs the changed factors and their
complete inputs, not a new calculation over every untouched place in the
system. This is the precise mathematical answer to the opening question of
how a local action can be assessed without a global recalculation each time.

Local evaluation does not license incomplete observation. If a factor couples
the source to a downstream or environmental coordinate, that entire factor
must be reevaluated. If an action changes an energy carrier as well as water,
the complete physical increment must say so. Computational locality follows
from a declared factorization; it is not a promise that the physical world
has no long-range dependencies.

\section{The Gaussian expression is a theorem below the definition}
For the quadratic potential of Chapter~\ref{ch:potential}, write
$H=\operatorname{diag}(1/\sigma_i^2)$ and $\mu(x^{\mathrm{b}})=H(x^{\mathrm{b}}-x^*)$. Exact expansion
of the squares gives
\begin{equation}\label{eq:gaussian-exact-ebu}
 E_e(q)=-q\,\mu(x^{\mathrm{b}})^T S_e-\frac{q^2}{2}S_e^T H S_e.
\end{equation}
The first term uses the initial slope. The second tells us how the slope
changes over the finite action. No small-$q$ approximation has been made.
For a lossless transfer from $i$ to $j$, $S_e=-e_i+e_j$, and the expression
becomes
\begin{equation}
 E_e(q)=q(\mu_i-\mu_j)
       -\frac{q^2}{2}\left(\frac1{\sigma_i^2}+\frac1{\sigma_j^2}\right).
\end{equation}
This edge polynomial is a corollary for this physical increment and this
quadratic potential. It is not the general definition of EBU. A lossy route,
a valued sink or a coupled potential changes the specialized formula while
leaving Equation~\ref{eq:exact-local-ebu} intact.

\section{One calculation with three independent readings}
Return to the lossless two-cell state $(18,2)$, reference $(10,10)$ and
scales $(2,2)$. A four-unit transfer reaches $(14,6)$. Direct evaluation
gives $V_{\mathrm{pre}}=16$ and $V_{\mathrm{post}}=4$, hence $E=12$.

The Gaussian expansion gives the same result: the initial edge field is
$4$ per stock unit and the directional curvature is $1/2$, so
$E=4(4)-(1/2)(4^2)/2=12$. Along the linear evaluation path the field falls
from $4$ to $2$; its average is $3$, and integrating over four units again
gives $12$. These are three derivations of one number, not three independent
claims of value.

The agreement is useful for checking software, but it does not establish that
four units were available, that a pump ran, or that a clinic was served.
Physical permission comes from the action model; occurrence comes from the
event record; service comes from a separately declared outcome condition.

\section{The cost of an omitted carrier}
If a pump consumes electricity and that consumption matters to the claim,
one may extend the physical state with an energy coordinate and declare the
pump's full transition. The potential over that extended state then evaluates
its chosen consequences. If some loss crosses the model boundary, the
boundary crossing must be explicit. By contrast, adding a number called
``process burden'' to the right of Equation~\ref{eq:exact-local-ebu} would
create a different account. It cannot silently redefine the physical EBU.

There is room for engineering cost ledgers and institutional charges. Their
units, provenance and decision role should be written alongside EBU, not
confused with it. That separation lets a reader ask which question each
number answers and prevents the same physical loss from being counted twice.

\section{From a quote to a record}
Before execution, a quote names a possible predecessor, potential version,
action increment and quantity. After execution, an event record must say
which quantity was actually accepted and which complete successor occurred.
The recorded EBU is recalculated from that pair. If the world changed between
quotation and execution, the old quote is an historical proposal, not the
current baseline.

The number can then be attributed to actors only under a declared rule. An
exact group total does not by itself decide who owns it or what anyone may
spend. We will first follow the physical event into its record, then return
to the more difficult case of several actions sharing one field.
